% Source: https://artofproblemsolving.com/community/c3h2023929_mathcon_2019_finals_grade_5
% Stable identity: aops-topic-2023929-q01
Moving from $a$ to $b$ on the number line covers a signed distance $b-a$. Going $\dfrac{2}{3}$ of the way means adding $\dfrac{2}{3}$ of this distance to the starting point:
\[
a+\dfrac{2}{3}(b-a)=\dfrac{1}{3}a+\dfrac{2}{3}b=\boxed{\dfrac{a+2b}{3}}.
\]
This is a weighted average of $a$ and $b$ whose weights $\dfrac{1}{3}$ and $\dfrac{2}{3}$ add to $1$, with more weight on $b$ because the point is closer to $b$.

The expression $\dfrac{2}{3}a+\dfrac{2}{3}b$ has weights adding to $\dfrac{4}{3}$, so it is not between $a$ and $b$ in general: for $a=b=3$ it gives $4$, but the point $\dfrac{2}{3}$ of the way from $3$ to $3$ is $3$. (For $a=0$, $b=3$ the correct answer is $2$, which $\dfrac{a+2b}{3}$ gives.)
